\subsection{Интегрирование простейших дробей 1--4 типов. Метод неопределённых коэффициентов}

\subsubsection*{1. Разложение правильной рациональной дроби}

Правильная рациональная дробь $\dfrac{P(x)}{Q(x)}$ (где $\deg P < \deg Q$) разлагается в сумму \textbf{простейших (парциальных) дробей} четырёх типов:

\begin{center}
\renewcommand{\arraystretch}{2.0}
\begin{tabular}{|c|l|l|}
\hline
\textbf{Тип} & \textbf{Вид дроби} & \textbf{Соответствует множителю $Q(x)$} \\
\hline
1 & $\dfrac{A}{x - a}$ & простой корень $a$ \\
2 & $\dfrac{A}{(x-a)^k}$, $k \ge 2$ & корень $a$ кратности $k$ \\
3 & $\dfrac{Ax + B}{x^2+px+q}$, $p^2-4q<0$ & неприводимый квадратный трёхчлен \\
4 & $\dfrac{Ax+B}{(x^2+px+q)^k}$, $k\ge 2$ & неприводимый трёхчлен кратности $k$ \\
\hline
\end{tabular}
\end{center}

\subsubsection*{2. Интегрирование простейших дробей}

\textbf{Тип 1:}
\[
\int \frac{A}{x-a}\,dx = A\ln|x-a| + C.
\]

\textbf{Тип 2:}
\[
\int \frac{A}{(x-a)^k}\,dx = \frac{A}{(1-k)(x-a)^{k-1}} + C, \quad k \ge 2.
\]

\textbf{Тип 3.} Знаменатель приводим к стандартной форме. Числитель разбивается на два слагаемых:
\[
\int \frac{Ax+B}{x^2+px+q}\,dx
= \frac{A}{2}\ln(x^2+px+q)
+ \frac{B - \tfrac{Ap}{2}}{\sqrt{q - \tfrac{p^2}{4}}}\,
\arctan\frac{x + \tfrac{p}{2}}{\sqrt{q - \tfrac{p^2}{4}}} + C.
\]

\textbf{Тип 4.} Вычисляется рекуррентно с помощью формулы понижения степени.

\subsubsection*{3. Метод неопределённых коэффициентов}

Неизвестные коэффициенты $A, B, \ldots$ в разложении находятся одним из двух способов:

\textbf{Способ 1 (приравнивание коэффициентов):} привести обе части к общему знаменателю и приравнять коэффициенты при одинаковых степенях $x$.

\textbf{Способ 2 (подстановка корней):} подставить в тождество значения $x$, являющиеся корнями знаменателя.

\subsubsection*{4. Пример: $\displaystyle\int \frac{dx}{x^3-1}$}

Разложим $x^3-1 = (x-1)(x^2+x+1)$.

\textbf{Шаг 1.} Запишем разложение на простейшие дроби:
\[
\frac{1}{x^3-1} = \frac{A}{x-1} + \frac{Bx+C}{x^2+x+1}.
\]

\textbf{Шаг 2.} Умножаем на $x^3-1$:
\[
1 = A(x^2+x+1) + (Bx+C)(x-1).
\]

\textbf{Шаг 3.} Находим коэффициенты:
\begin{itemize}
    \item $x=1$: $1 = 3A \implies A = \dfrac{1}{3}$.
    \item Приравниваем коэффициенты при $x^2$: $0 = A+B \implies B = -\dfrac{1}{3}$.
    \item Приравниваем свободные члены: $1 = A - C \implies C = A-1 = -\dfrac{2}{3}$.
\end{itemize}

\textbf{Шаг 4.} Интегрируем:
\[
\int \frac{dx}{x^3-1}
= \frac{1}{3}\int\frac{dx}{x-1}
+ \int\frac{-\tfrac{1}{3}x - \tfrac{2}{3}}{x^2+x+1}\,dx.
\]

Второй интеграл разбиваем:
\[
\int\frac{-\tfrac{1}{3}x-\tfrac{2}{3}}{x^2+x+1}\,dx
= -\frac{1}{6}\int\frac{2x+1}{x^2+x+1}\,dx
- \frac{1}{2}\int\frac{dx}{x^2+x+1}.
\]

Первое слагаемое: $-\dfrac{1}{6}\ln(x^2+x+1)$.

Второе: приводим $x^2+x+1 = \left(x+\tfrac{1}{2}\right)^2 + \dfrac{3}{4}$,
\[
-\frac{1}{2}\int\frac{dx}{\left(x+\frac{1}{2}\right)^2+\frac{3}{4}}
= -\frac{1}{2}\cdot\frac{2}{\sqrt{3}}\arctan\frac{x+\frac{1}{2}}{\frac{\sqrt{3}}{2}}
= -\frac{1}{\sqrt{3}}\arctan\frac{2x+1}{\sqrt{3}}.
\]

\textbf{Ответ:}
\[
\boxed{
\int\frac{dx}{x^3-1}
= \frac{1}{3}\ln|x-1|
- \frac{1}{6}\ln(x^2+x+1)
- \frac{1}{\sqrt{3}}\arctan\frac{2x+1}{\sqrt{3}} + C.
}
\]